\documentclass[11pt,a4paper]{article}
\usepackage[margin=25mm]{geometry}
\usepackage{amsmath,amssymb,amsthm,hyperref}
\setlength{\emergencystretch}{2em}
\newtheorem{theorem}{Theorem}
\newtheorem{corollary}[theorem]{Corollary}
\newcommand{\gen}[1]{\langle#1\rangle}
\title{Orbit deficits on a contiguous interval}
\author{Human Compression project: L044 continuation}
\date{}
\begin{document}\maketitle
Let $G=\gen{\rho_0,\ldots,\rho_{r-1}}$ act transitively on a
finite set $\Omega$ of size $n$, where the generators are involutions
and $[\rho_j,\rho_k]=1$ for $|j-k|>1$. An edge of colour $j$
is \emph{ambient-eligible} if its endpoints lie in different orbits
of $G_j=\gen{\rho_k:k\ne j}$.
Fix a nonempty contiguous interval $I=\{a,\ldots,b\}$, set
$H=\gen{\rho_j:j\in I}$ and $U=\operatorname{supp}(H)$, and
suppose that every colour in $I$ has at least two ambient-eligible
edges. Select exactly two per colour on vertex set $U$.

\begin{theorem}\label{interval}
In a square-only selection minimizing the number of squares,
every proper nontrivial $H$-orbit has a tree component. If $H$
is intransitive and has $c$ nontrivial orbits and $o(H)$ total
orbits on $\Omega$, then
\[
  2|I|\leq |U|-c=n-o(H).                              \tag{1}
\]
No intersection property, parity assumption or primitivity is needed.
\end{theorem}
\begin{proof}
We use the complete accepted proof of the ambient terminal-orbit
deficit, including its paid paired-cut and equality-localisation
proofs. Here is the exact adaptation; none of those arguments is
being applied with local $H_j$ eligibility.

The orbit-cut exchanges and square-block lemma remain valid for
ambient-eligible edges of any selected colours. The square-only
normal form and its minimum-square property use ordinary slides
and the global middle move. When a square has colours $i-1,i+1$
in $I$, contiguity gives $a<i<b$, so $\rho_i\in H$.
Thus these moves preserve actual $H$-orbits. In particular every
minimum-square selection has at most one square per component,
just as in the terminal proof.

Fix a proper nontrivial $H$-orbit $O$ and suppose it has no tree
component. Its selected edge count is then $E_O=|O|$.
The localisation lemma, the ordinary-slide partition argument and
the middle-transport argument from the accepted proof are unchanged:
all square-colour eligible edges lie in the square's component;
ordinary slides preserve its vertex set; and middle moves transport
the component partition in $O$ while preserving $E_O=|O|$.
These statements used only $\rho_i\in H$, not that $I$ was terminal.
In particular every full-graph neighbour of a square vertex stays
in that selected component.

Full transitivity supplies a full edge leaving $O$. Minimize, over
all minimum-square selections having $E_O=|O|$, a connector consisting
of a selected path from a square to the inside endpoint of such an
edge, followed by that single exit edge. Its exit colour $x$ is
outside $I$, because each generator in $I$ preserves $O$.
The selected path is nonempty by the neighbour assertion.
An ordinary first-edge slide shortens the connector by the accepted
localisation argument, including its exclusion of a second occurrence
of the first colour farther down the selected path.

In the remaining case the first edge has colour $i$ and the square
colours are $i-1,i+1$. Since $a<i<b$ and $x\notin[a,b]$, we have
$|i-x|\geq2$. Hence $\rho_i$ commutes with the exit generator
on \emph{either} side of the interval. Retain the square and transport
all other selected edges by $\rho_i$. Transport the connector after
its first edge as well, including its full exit edge. This is an
actual exit edge of colour $x$; its outside endpoint remains outside
$O$, since $\rho_i\in H$. Its selected path now begins on the retained
square and is shorter by one. Any coincidence can be removed from
the walk to shorten it further. This is the same contradiction as
in the terminal proof. Therefore $O$ has a tree component.

For intransitive $H$, every nontrivial orbit is proper. Each supplies
at least one unit of vertices minus edges; square components supply
zero. Summing gives $2|I|\leq |U|-c$. The fixed points contribute
$n-|U|$ further $H$-orbits, giving the final equality in (1).
\end{proof}

\begin{corollary}[A middle interval with two outside budgets]
Suppose $0\leq h<g\leq r-1$, $g\geq h+2$, and every colour
in $I=\{h+1,\ldots,g-1\}$ has two ambient-eligible edges.
Let $H=G_I$ be intransitive, with support $V$ and $c$ nontrivial
orbits. Suppose sets $X,Y\subseteq\Omega$ satisfy
\[
 V\subseteq X\cap Y,\quad X\cup Y=\Omega,\quad
 2h\leq n-|X|-1,\quad 2(r-1-g)\leq n-|Y|-1.
\]
Then $2r\leq n+2-c$. In particular $c\geq3$ implies
$r\leq\lfloor(n-1)/2\rfloor$. If $c=2$, the only failure
of that target has even $n$ and $2r=n$, and forces equality
in both outside budgets, $V=X\cap Y$, and exactly one tree
component in each nontrivial $H$-orbit of a minimum selection.
\end{corollary}
\begin{proof}
Theorem~\ref{interval} and the two budgets give
\[
\begin{split}
2r&=2h+2(g-h-1)+2(r-1-g)+4\\
&\leq(n-|X|-1)+(|V|-c)+(n-|Y|-1)+4\\
&\leq n+2-c.
\end{split}
\]
For $c=2$ every inequality must be an equality if the target fails.
Equality in the orbit deficit says that the total number of tree
components is $c$, proving the stated necessary conditions.
\end{proof}

\paragraph{Source routing and attribution.}
This extends the accepted ambient version of the rank step in
Cameron--Fernandes--Leemans--Mixer, \emph{Highest rank of a polytope
for $A_n$}, arXiv:1605.09173v1, Proposition~5.27. It is supporting
theory for auditing Proposition~5.28, not a recovery of that
proposition: the stated support and outside-budget hypotheses of the
corollary must still be proved in an application. In particular the
reversal of a prefix index is $r-1-g$, and the two-orbit equality
case is not discarded. Local $H_j$ eligibility is not substituted
for ambient eligibility.

The full three-page terminal proof and its accepted paired-cut and
equality proofs remain paid dependencies. The earlier global
middle-colour move was supplied by the root collaborator and received
a distinct review; the two-sided exit adaptation here is the present
author's continuation. No census, new formal theorem, whole L044
recovery, original credit or question credit is claimed.
\end{document}
